\begin{figure}[htbp]
\centering
\begin{tikzpicture}

% ============================================================================
% Panel (a): Impulse function delta(t)
% ============================================================================
\begin{axis}[
    name=impulse,
    width=5.5cm, height=4.5cm,
    xlabel={$t$},
    ylabel={$\delta(t)$},
    xmin=-1.5, xmax=3,
    ymin=-0.2, ymax=1.5,
    xtick={-1, 0, 1, 2},
    ytick={0, 1},
    grid=major,
    grid style={gray!30},
    title={\small (a) Impulse},
    title style={at={(0.5,1.02)}, anchor=south},
    axis lines=middle,
    axis line style={->, >=stealth, thick},
    xlabel style={at={(axis description cs:1.05,0.15)}},
    ylabel style={at={(axis description cs:0.15,1.05)}}
]

% Zero line for t < 0 and t > 0
\addplot[UofSCGarnet, very thick, domain=-1.5:-0.01, samples=50] {0};
\addplot[UofSCGarnet, very thick, domain=0.01:3, samples=50] {0};

% Impulse arrow at t = 0
\draw[UofSCGarnet, very thick, ->, >=stealth] (axis cs:0,0) -- (axis cs:0,1.2);
\node[UofSCGarnet, font=\footnotesize, right] at (axis cs:0.1,1.1) {$\infty$};

% Filled square marker at base
\node[rectangle, fill=UofSCGarnet, minimum size=4pt, inner sep=0pt] at (axis cs:0,0) {};

\end{axis}

% ============================================================================
% Panel (b): Unit step function u(t)
% ============================================================================
\begin{axis}[
    name=step,
    at={(impulse.outer east)},
    anchor=outer west,
    xshift=0.8cm,
    width=5.5cm, height=4.5cm,
    xlabel={$t$},
    ylabel={$u(t)$},
    xmin=-1.5, xmax=3,
    ymin=-0.2, ymax=1.5,
    xtick={-1, 0, 1, 2},
    ytick={0, 1},
    grid=major,
    grid style={gray!30},
    title={\small (b) Unit Step},
    title style={at={(0.5,1.02)}, anchor=south},
    axis lines=middle,
    axis line style={->, >=stealth, thick},
    xlabel style={at={(axis description cs:1.05,0.15)}},
    ylabel style={at={(axis description cs:0.15,1.05)}}
]

% Step function: 0 for t < 0, 1 for t >= 0
\addplot[UofSCAtlantic, very thick, domain=-1.5:-0.01, samples=50] {0};
\addplot[UofSCAtlantic, very thick, domain=0:3, samples=50] {1};

% Vertical line at discontinuity
\draw[UofSCAtlantic, very thick] (axis cs:0,0) -- (axis cs:0,1);

% Markers: open circle at (0,0), filled at (0,1)
\node[circle, draw=UofSCAtlantic, fill=white, minimum size=4pt, inner sep=0pt, line width=0.8pt] at (axis cs:0,0) {};
\node[circle, fill=UofSCAtlantic, minimum size=4pt, inner sep=0pt] at (axis cs:0,1) {};

\end{axis}

% ============================================================================
% Panel (c): Ramp function r(t) = t*u(t)
% ============================================================================
\begin{axis}[
    name=ramp,
    at={(step.outer east)},
    anchor=outer west,
    xshift=0.8cm,
    width=5.5cm, height=4.5cm,
    xlabel={$t$},
    ylabel={$r(t)$},
    xmin=-1.5, xmax=3,
    ymin=-0.2, ymax=2.5,
    xtick={-1, 0, 1, 2},
    ytick={0, 1, 2},
    grid=major,
    grid style={gray!30},
    title={\small (c) Ramp},
    title style={at={(0.5,1.02)}, anchor=south},
    axis lines=middle,
    axis line style={->, >=stealth, thick},
    xlabel style={at={(axis description cs:1.05,0.1)}},
    ylabel style={at={(axis description cs:0.15,1.05)}}
]

% Ramp function: 0 for t < 0, t for t >= 0
\addplot[UofSCHorseshoe, very thick, domain=-1.5:0, samples=50] {0};
\addplot[UofSCHorseshoe, very thick, domain=0:2.4, samples=50] {x};

% Corner marker
\node[circle, fill=UofSCHorseshoe, minimum size=4pt, inner sep=0pt] at (axis cs:0,0) {};

% Slope annotation
\node[UofSCHorseshoe, font=\footnotesize] at (axis cs:1.6,0.6) {slope $= 1$};

\end{axis}

% ============================================================================
% Panel (d): Parabolic function t^2*u(t)
% ============================================================================
\begin{axis}[
    name=parabola,
    at={(impulse.outer south west)},
    anchor=outer north west,
    yshift=-0.8cm,
    width=5.5cm, height=4.5cm,
    xlabel={$t$},
    ylabel={$\frac{1}{2}t^2 u(t)$},
    xmin=-1.5, xmax=3,
    ymin=-0.3, ymax=4,
    xtick={-1, 0, 1, 2},
    ytick={0, 1, 2, 3},
    grid=major,
    grid style={gray!30},
    title={\small (d) Parabolic},
    title style={at={(0.5,1.02)}, anchor=south},
    axis lines=middle,
    axis line style={->, >=stealth, thick},
    xlabel style={at={(axis description cs:1.05,0.1)}},
    ylabel style={at={(axis description cs:0.22,1.05)}}
]

% Parabolic function: 0 for t < 0, 0.5*t^2 for t >= 0
\addplot[UofSCCongaree, very thick, domain=-1.5:0, samples=50] {0};
\addplot[UofSCCongaree, very thick, domain=0:2.7, samples=100] {0.5*x^2};

% Corner marker
\node[circle, fill=UofSCCongaree, minimum size=4pt, inner sep=0pt] at (axis cs:0,0) {};

\end{axis}

% ============================================================================
% Panel (e): Sinusoidal function sin(omega*t)
% ============================================================================
\begin{axis}[
    name=sinusoid,
    at={(parabola.outer east)},
    anchor=outer west,
    xshift=0.8cm,
    width=11.6cm, height=4.5cm,
    xlabel={$t$},
    ylabel={$\sin(\omega t)$},
    xmin=-0.5, xmax=8,
    ymin=-1.4, ymax=1.4,
    xtick={0, 1.571, 3.142, 4.712, 6.283},
    xticklabels={$0$, $\frac{T}{4}$, $\frac{T}{2}$, $\frac{3T}{4}$, $T$},
    ytick={-1, 0, 1},
    yticklabels={$-A$, $0$, $A$},
    grid=major,
    grid style={gray!30},
    title={\small (e) Sinusoid},
    title style={at={(0.5,1.02)}, anchor=south},
    axis lines=middle,
    axis line style={->, >=stealth, thick},
    xlabel style={at={(axis description cs:1.02,0.45)}},
    ylabel style={at={(axis description cs:0.05,1.05)}}
]

% Sinusoidal function
\addplot[UofSCRose, very thick, domain=-0.5:8, samples=200] {sin(deg(x))};

% Period annotation
\draw[UofSC70Black, <->, >=stealth, thick] (axis cs:0,-1.25) -- (axis cs:6.283,-1.25);
\node[UofSC70Black, font=\footnotesize, below] at (axis cs:3.14,-1.25) {Period $T = \frac{2\pi}{\omega}$};

% Amplitude annotation
\draw[UofSCRose, <->, >=stealth] (axis cs:7.2,-1) -- (axis cs:7.2,1);
\node[UofSCRose, font=\footnotesize, right] at (axis cs:7.2,0) {$2A$};

\end{axis}

\end{tikzpicture}
\caption{Standard test input signals for control systems analysis. (a) \textbf{Impulse function} $\delta(t)$: infinite at $t=0$, zero elsewhere, with unit area---used to characterize system impulse response. (b) \textbf{Unit step function} $u(t)$: discontinuous jump from 0 to 1 at $t=0$---the most common test signal for transient response. (c) \textbf{Ramp function} $r(t) = tu(t)$: linearly increasing signal---tests steady-state tracking of changing references. (d) \textbf{Parabolic function} $\frac{1}{2}t^2 u(t)$: quadratic growth---tests acceleration tracking capability. (e) \textbf{Sinusoidal function} $A\sin(\omega t)$: periodic oscillation with amplitude $A$ and period $T = 2\pi/\omega$---fundamental for frequency response analysis.}
\label{fig:test_signals}
\end{figure}
